\title{FlashAttention-3: Fast and Accurate Attention with Asynchrony and Low-precision}

\begin{abstract}
The attention mechanism, fundamental to the Transformer model widely used in modern natural language processing, constitutes a major computational bottleneck for scaling large language models and extending context length. 
\textsc{FlashAttention} introduced an algorithmic strategy for accelerating attention on GPUs by optimizing memory access. Nevertheless, this approach does not fully utilize advanced functionalities introduced in latest hardware generations, with \textsc{FlashAttention-2} attaining merely 35\% computational efficiency on the H100 GPU. In this work, we propose three key optimizations for enhancing attention throughput on Hopper GPUs: leveraging the inherent asynchrony between Tensor Cores and TMA to (1) maximize overlap of computation and data transfers using warp-specialization, (2) perform simultaneous block-wise matrix multiplication and softmax, and (3) employ block quantization together with incoherent execution, exploiting dedicated FP8 precision support in hardware. Empirical results demonstrate that our proposed \textsc{FlashAttention-3} offers 1.5-2.0$\times$ acceleration on H100 GPUs, achieving up to 740 TFLOPs/s with FP16 (corresponding to 75\% utilization) and approaching 1.2 PFLOPs/s with FP8 precision. Furthermore, FP8 \textsc{FlashAttention-3} yields 2.6$\times$ reduced numerical error compared to a standard FP8 attention implementation.
\end{abstract}

\section{Introduction}
\label{sec:intro}

Within the Transformer architecture~\citep{vaswani2017attention}, the attention mechanism emerges as the main computational constraint due to the self-attention score calculations between queries and keys, which scale quadratically with the sequence length. Extending attention mechanisms to handle substantially longer contexts has the potential to enable novel capabilities, such as complex reasoning across multiple extended documents~\citep{guo2021longt5,shaham2022scrolls,peng2023yarn}, navigating large software codebases~\citep{roziere2023code, li2023starcoder}, processing new data types such as high-resolution imagery~\citep{chen2022scaling}, audio sequences~\citep{gulati2020conformer}, and video~\citep{ho2022video}, along with supporting broader applications like user engagement histories~\citep{sun2019bert4rec} and long-term planning for agents~\citep{yao2022react}. As a result, there has been a surge in research on accelerating attention for lengthy contexts, which includes methods based on approximation~\citep{katharopoulos2020transformers,choromanski2020rethinking,
tay2020efficient}, advancements in software-level optimizations (\citep{rabe2021self,dao2022flashattention,kwon2023efficient}), and the exploration of alternative architectural frameworks~\citep{peng2023rwkv,sun2023retentive,gu2023mamba}.

This paper extends the approach introduced by \citet{dao2022flashattention}, which designed precise attention algorithms by factoring GPU architectural features and execution patterns into algorithmic structure. In their seminal work, \citep{dao2022flashattention} presented \textsc{FlashAttention}, a technique employing a unique tiling method for parallelizing attention computations. By consolidating all attention-related operations into a single GPU kernel, this method avoids unnecessary accesses to slow global memory.

Later, \citet{dao2023flashattention2} reformulated the algorithm as \textsc{FlashAttention-2}, which further exploits parallelism across the sequence length dimension. Additionally, it executes the forward pass’s inner loop over blocks of the key and value matrices, optimizing workload distribution and occupancy on the GPU. Nevertheless, our analysis reveals that \textsc{FlashAttention-2} continues to underperform regarding utilization on advanced GPUs, especially when compared to highly optimized matrix-multiplication (GEMM) kernels; for instance, it achieves only 35\% utilization compared to 80-90\% on the Hopper H100 GPU. This discrepancy is partly due to differences at the implementation level, including reliance on Ampere instructions rather than leveraging Hopper-specific instructions for Tensor Core optimization. Prior efforts, such as ThunkerKitten~\citep{spector2024thunder} and cuDNN 9~\citep{cudnn9}, have demonstrated that attention can be made faster and code complexity reduced by adopting Hopper-optimized instructions and tile-based abstractions.

At a foundational level, the algorithm underlying \textsc{FlashAttention-2} operates within a streamlined synchronous framework, lacking integration of asynchronous execution and low-precision arithmetic. Asynchrony emerges from specialized hardware tailored for the acceleration of critical machine learning operations: dedicated units for matrix multiplication (Tensor Cores) and memory access (Tensor Memory Accelerator -- TMA) functioning independently from CUDA cores that handle logical, integer, and floating-point operations. The adoption of lower numerical precision—such as FP8 in Hopper and FP4 in Blackwell, continuing with FP16 (Pascal, 2017) and BF16 (Ampere, 2020)—has consistently enabled increased computational throughput (twofold or fourfold) without additional demands on power consumption or chip space. An overview of Hopper's advancements in these aspects is presented in \cref{subsec:hardware}. 

The principal engineering challenge involves revising \textsc{FlashAttention-2} so it can leverage these advanced hardware capabilities. Incorporating asynchrony necessitates orchestrating the concurrent execution of matrix multiplication and softmax operations despite inherent data dependencies, while introducing low-precision techniques requires diligent efforts to curb quantization artifacts, a concern heightened by atypical or outlier values frequently encountered in LLMs~\citep{dettmers2208llm, sun2024massive}.

With this objective in mind, we introduce \textsc{FlashAttention-3}, which integrates three innovative approaches to enhance efficiency on the latest GPU platforms:\footnote{Our findings are presented with respect to NVIDIA's Hopper architecture. Nonetheless, the algorithm is applicable to any GPU architecture that supports strong asynchronous operations and low-precision computation.}

\begin{enumerate}

\item \textbf{Producer-Consumer asynchrony:} We introduce a warp-specialized software pipelining method that leverages the asynchronous nature of data transfer and Tensor Core computation. By assigning distinct warps to the roles of data producers and consumers, this approach enhances the pipeline's capacity to mask memory access delays and instruction scheduling overheads.
\item \textbf{Hiding softmax under asynchronous block-wise GEMMs:} The algorithm is modified to allow non-GEMM softmax-related operations, including floating point multiply-add and exponentiation—which are lower in throughput—to occur concurrently with block-level GEMM instructions issued via asynchronous WGMMA. Specifically, we restructure the \textsc{FlashAttention-2} workflow to avoid certain sequential interlocks between softmax and GEMM computation. As an illustrative case: within the algorithm's two-stage arrangement, the softmax operates on one portion of the score matrix while the asynchronous WGMMA computes the subsequent block in parallel.
\item \textbf{Hardware-accelerated low-precision GEMM:} The forward algorithm is reconfigured to utilize FP8 Tensor Cores for GEMM computation, resulting in an approximate doubling of observed TFLOPs/s. Addressing the challenge of disparate memory layouts required by WGMMA for FP32 accumulators and FP8 operands, we employ block quantization and incoherent processing strategies, thereby curtailing the reduction in accuracy associated with transitioning to FP8 precision.
\end{enumerate}

For empirical validation, we conduct benchmarking of \textsc{FlashAttention-3} on the H100 SXM5 GPU, varying several parameters. Our results indicate: (1) FP16 provides a speedup of 1.5 to 2.0$\times$ in the forward pass compared to \textsc{FlashAttention-2} (with peak throughput up to 740 TFLOPs/s) and accelerates the backward pass by 1.5 to 1.75$\times$, (2) FP8 achieves nearly 1.2 PFLOPs/s performance, and (3) at large sequence lengths, FP16 exhibits superior performance, while FP8 remains competitive\footnote{Specifically, for head dimension 64, FP8 in \textsc{FlashAttention-3} surpasses other approaches; for head dimensions of 128 and 256, FP8 performs comparably in non-causal scenarios but is outperformed with causal masking.} against NVIDIA's cuDNN library's leading attention implementation. Furthermore, we confirm that FP16 in \textsc{FlashAttention-3} produces the same level of numerical error as \textsc{FlashAttention-2}, and improves upon standard attention due to retaining intermediate computations (such as softmax scaling) in FP32 precision. Additionally, for cases featuring outlier elements, FP8 in \textsc{FlashAttention-3} using block quantization and incoherent processing yields 2.6$\times$ higher accuracy compared to standard attention implementations that utilize per-tensor quantization.

\textsc{FlashAttention-3} has been released under an open and permissive license\footnote{\textsc{FlashAttention-3} can be found at \texttt{https://github.com/Dao-AILab/flash-attention}}. We aim to incorporate it into both PyTorch and Hugging Face libraries, maximizing its accessibility to the research and developer communities.

\section{Background: Multi-Head Attention and GPU Characteristics}
\label{sec:background}

\subsection{Multi-Head Attention}
\label{subsec:multi_head_attn}

Let $\mathbf{Q}, \mathbf{K}, \mathbf{V} \in \mathbb{R}^{N \times d}$ be the query, key and value input sequences associated to a single head, where $N$ is the sequence length and $d$ is the head dimension. Then the attention output $\mathbf{O}$ is computed as:

\begin{equation*}
  \mathbf{S} = \alpha \mathbf{Q} \mathbf{K}^\top \in \mathbb{R}^{N \times N}, \quad \mathbf{P} = \mathrm{softmax}(\mathbf{S}) \in \mathbb{R}^{N \times N}, \quad \mathbf{O} = \mathbf{P}\mathbf{V} \in \mathbb{R}^{N \times d},
\end{equation*}

Here, the $\mathrm{softmax}$ function operates across each row, and it is standard to choose $\alpha = 1/\sqrt{d}$ as the normalization constant. To mitigate numerical issues arising from exponentiation, it is common to subtract $\mathrm{rowmax}(\mathbf{S})$ from $\mathbf{S}$. In the multi-head attention (MHA) framework, each individual head is equipped with separate projections for queries, keys, and values; this allows for simultaneous computation across different heads and batches, ultimately generating the complete output tensor.

Now let $\phi$ be a scalar loss function and let $\mathbf{d}(-) = \partial \phi / \partial (-)$ be notation for the gradient. Given the output gradient $\mathbf{dO} \in \mathbb{R}^{N \times d}$, we compute $\mathbf{dQ}$, $\mathbf{dK}$, and $\mathbf{dV}$ according to the chain rule as follows:

\begin{align*}
  \mathbf{dV} &= \mathbf{P}^\top \mathbf{dO} \in \mathbb{R}^{N \times d} \\
  \mathbf{dP} &= \mathbf{dO} \mathbf{V}^\top \in \mathbb{R}^{N \times N} \\
  \mathbf{dS} &= \mathrm{dsoftmax} (\mathbf{dP}) \in \mathbb{R}^{N \times N} \\
  \mathbf{dQ} &= \alpha \mathbf{dS} \mathbf{K} \in \mathbb{R}^{N \times d} \\
  \mathbf{dK} &= \alpha \mathbf{dS}^\top \mathbf{Q} \in \mathbb{R}^{N \times d},
\end{align*}

In this context, the relationship $\mathbf{d}s = (\mathrm{diag}(p) - p p^\top)\mathbf{d}p$ holds, where $p$ is defined as $p = \mathrm{softmax}(s)$ with $s$ being a vector. We denote $\mathrm{dsoftmax}(\mathbf{dP})$ to represent this operation executed on each row individually. This calculation, as with the previous step, is efficiently parallelized over the heads and batches when performing the backward propagation in MHA.

\subsection{GPU hardware characteristics and execution model}
\label{subsec:hardware}

This section outlines features of the GPU execution model pertinent to \textsc{FlashAttention-3}, emphasizing the NVIDIA Hopper architecture as a specific example of this framework.

\paragraph{Memory hierarchy:} The organization of GPU memory follows a tiered structure, where larger storage capacities are typically associated with lower bandwidths (\cref{tab:gpu-hierarchy})\footnote{\citet{luo2024benchmarking} estimates the shared memory bandwidth as 128 bytes per clock cycle per SM; for total throughput, this is multiplied by 132 SMs and the maximum clock rate of 1830 MHz.}. At the base level, global memory (GMEM)—or HBM—serves as the external DRAM, which all streaming multiprocessors (SMs) can access. Data residing in GMEM can be cached seamlessly in an on-chip L2 cache. Progressing down the hierarchy, each SM is equipped with a limited amount of fast, multi-banked shared memory (SMEM), which must be explicitly managed by the programmer. The final, fastest level consists of the register files, located within each SM.

\paragraph{Thread hierarchy:} Within the GPU programming model, execution is structured using nested collections of threads. The hierarchy starts at the individual thread and aggregates upwards: threads are organized into warps consisting of 32 threads, with four consecutive warps forming a warpgroup. These warpgroups are in turn grouped into threadblocks (also referred to as cooperative thread arrays or CTAs). In the Hopper architecture, multiple threadblocks can be assembled into threadblock clusters, and at the highest level, grids encompass all of these constructs.

The interconnection between these hierarchies is significant. All threads belonging to a single CTA are executed concurrently on the same SM, while CTAs grouped within a cluster share scheduling on the same GPC. SMEM is accessible by every thread inside a CTA, but each thread individually possesses up to 256 dedicated registers (RMEM) that are not shared with others.

\paragraph{Asynchrony and warp-specialization:}

GPUs operate as throughput-oriented processors, utilizing parallelism and asynchronous execution to mitigate delays associated with memory access and computation. In the context of asynchronous memory transfers between GMEM and SMEM, Hopper features the Tensor Memory Accelerator (TMA), a specialized hardware component \cite[\S7.29]{cuda}. Additionally, Hopper’s Tensor Core, accessible through the warpgroup-level WGMMA instruction \cite[\S9.7.14]{ptx}, differs from earlier architectures like Ampere by functioning asynchronously and enabling direct input sourcing from shared memory.

Support for asynchrony at the hardware level enables the use of warp-specialized kernels, wherein each warp within a CTA assumes a dedicated role as either a producer or consumer, issuing exclusively data transfer or computational instructions. This separation facilitates the compiler’s capacity to construct more efficient instruction schedules \citep{warp-specialization-2011}. Furthermore, Hopper provides mechanisms for adjusting register allocation dynamically across warpgroups through \verb|setmaxnreg| \cite[\S9.7.17.1]{ptx}, which permits warps performing MMAs to access a greater portion of RMEM compared to warps restricted to executing TMA, tasks that require only a single thread.

\paragraph{Low-precision number formats:}
\label{sec:low-precision-gpu}
Contemporary GPUs feature dedicated hardware blocks optimized for low-precision calculations. As an illustration, the WGMMA instruction takes advantage of Hopper's FP8 Tensor Cores, providing twice the throughput per SM relative to FP16 or BF16.

Successful use of FP8 WGMMA requires awareness of the operand layout requirements. Consider a GEMM operation computing $A \times B^{\top}$, with $A$ as an $M \times K$ matrix and $B$ as an $N \times K$ matrix. The $A$ or $B$ operand is described as \emph{mn-major} when its elements are stored contiguously along the outer $M$ or $N$ dimensions, and as \emph{k-major} when the contiguous storage occurs in the inner $K$ dimension. FP16 WGMMA accommodates both mn-major and k-major input layouts for matrices in SMEM. In contrast, FP8 WGMMA restricts supported operand layouts to the k-major form only. This constraint creates difficulties in contexts like attention mechanisms, where fusing sequential GEMM operations within one kernel leads to mismatches between FP32 accumulator layouts and FP8 operand layouts, thereby complicating calls to dependent FP8 WGMMAs.

When applied to attention mechanisms, these layout constraints require specific adaptations in the structure of an FP8 algorithm, as detailed in \cref{sec:algofp8}.

\subsection{Standard Attention and Flash Attention} In accordance with \citet{dao2022flashattention}, we use the term \textbf{standard attention} to refer to the typical GPU-based attention computation, in which the intermediate matrices $\mathbf{S}$ and $\mathbf{P}$ are explicitly stored in HBM. The core principle behind \textsc{FlashAttention} is the exploitation of a localized softmax reduction, which circumvents the costly intermediate memory operations by fusing the attention steps into a unified computational kernel. This local softmax operation is captured in lines \ref{code-ws:softmax_start}-\ref{code-ws:softmax_end} within the mainloop of the consumer, as detailed in \cref{alg:flash3_wgmma_ws_only}, in conjunction with the blockwise rescaling of $\mathbf{O}$. The straightforward derivation confirming that this approach indeed yields the intended computation of $\mathbf{O}$ is provided in \cite[\S 2.3.1]{dao2023flashattention2}.

\section{FlashAttention-3: Algorithm}
\label{sec:algo}

In this section, we present the \textsc{FlashAttention-3} algorithm. For clarity, our discussion centers on the forward pass, while the algorithm for the backward pass is outlined in~\cref{sec:algo_ws_bwd}. We begin by detailing how warp-specialization, combined with a circular SMEM buffer, can be incorporated into the core structure of \textsc{FlashAttention-2}. Next, we illustrate how leveraging the asynchronous nature of WGMMA allows the implementation of a two-stage pipeline that overlaps GEMM and softmax computations. Lastly, we outline the adjustments necessary for FP8, addressing both layout compatibility and enhanced numerical fidelity through block quantization and incoherent processing.

\subsection{Producer-Consumer asynchrony through warp-specialization and pingpong scheduling}
\label{sec:algo_ws}

\paragraph{Warp-specialization}
Similar to \textsc{FlashAttention-2}, \textsc{FlashAttention-3}'s forward computation exhibits significant parallelism across the batch dimension, the number of heads, and the query sequence axis. Therefore, providing an overview at the CTA (Cooperative Thread Array) granularity is adequate, wherein the algorithm processes a tile $\mathbf{Q}_i$ of the queries to produce the associated output tile $\mathbf{O}_i$. For clarity, we introduce the warp-specialization mechanism that utilizes a circular SMEM buffer, omitting for now any GEMM-softmax overlap. Denote $d$ as the head dimension and $N$ as the sequence length. By choosing a query block size $B_r$, the query matrix $\mathbf{Q}$ is partitioned into $T_r = \lceil \frac{N}{B_r} \rceil$ tiles $\mathbf{Q}_1, \ldots, \mathbf{Q}_{T_r}$.

\begin{algorithm}[H]
    \caption{\small\label{alg:flash3_wgmma_ws_only}\textsc{FlashAttention-3} forward computation \textbf{excluding} intra-consumer overlap -- CTA representation}
    \begin{algorithmic}[1]
\REQUIRE Input: matrices $\mathbf{Q}_i \in \mathbb{R}^{B_r \times d}$ and $\mathbf{K}, \mathbf{V} \in \mathbb{R}^{N \times d}$ resident in HBM, key block size $B_c$ where $T_c = \lceil \frac{N}{B_c} \rceil$.
\STATE Set up pipeline instance for synchronizing barriers with a $s$-stage circular SMEM buffering scheme.
\IF {executing within a producer warpgroup}
\STATE Release a preset quantity of registers.
\STATE Trigger data transfer of $\mathbf{Q}_i$ from HBM to SMEM.
\STATE When transfer finalizes, update status to inform consumers that $\mathbf{Q}_i$ is available.
\FOR{$0 \le j < T_c$}
    \STATE Block until consumption of the buffer stage indexed $(j\,\%\,s)$ completes.
    \STATE Trigger transfers of $\mathbf{K}_j, \mathbf{V}_j$ from HBM into SMEM at position $(j\,\%\,s)$ of the buffer.
    \STATE After transfers finish, mark the buffer stage as ready for consumption by consumers.
\ENDFOR
\ELSE \STATE Assign back registers proportionally to the active consumer warp count.
\STATE On the chip, set $\mathbf{O}_i = (0) \in \mathbb{R}^{B_r \times d}$ and assign initial values $\ell_i, m_i = (0), (-\infty) \in \mathbb{R}^{B_r}$.
\STATE Hold execution until $\mathbf{Q}_i$ is ready in SMEM.
\FOR{$0 \le j < T_c$}
\STATE Wait for $\mathbf{K}_j$ to reside in shared memory.
\STATE Evaluate $\mathbf{S}_i^{(j)} = \mathbf{Q}_i \mathbf{K}_j^T$ via SS-GEMM operation, then commit and await progression. \label{alg:ws_only_gemm1}
\STATE Save $m_i^{\mathrm{old}} = m_i$ and update $m_i = \mathrm{max}(m_i^{\mathrm{old}}, \mathrm{rowmax}(\mathbf{S}_i^{(j)}))$. \label{code-ws:softmax_start}
\STATE Compute $\widetilde{\mathbf{P}}_i^{(j)} = \mathrm{exp}(\mathbf{S}_i^{(j)} - m_i)$ and update $\ell_i = \mathrm{exp}(m_i^{\mathrm{old}} - m_i) \ell_i + \mathrm{rowsum}(\widetilde{\mathbf{P}}_i^{(j)})$. \label{code-ws:softmax_end}
\STATE Block until $\mathbf{V}_j$ resides in SMEM.
\STATE Update $\mathbf{O}_i = \mathrm{diag}(\mathrm{exp}(m_i^{\mathrm{old}} - m_i))^{-1} \mathbf{O}_i + \widetilde{\mathbf{P}}_i^{(j)} \mathbf{V}_j$ using RS-GEMM; mark commit and await. \label{alg:ws_only_gemm2}
\STATE Release buffer stage $(j\,\%\,s)$, enabling the producer to continue.
\ENDFOR
\STATE Perform final update: $\mathbf{O}_i = \mathrm{diag}(\ell_i)^{-1} \mathbf{O}_i$ and $L_i = m_i + \log(\ell_i)$.
\STATE Output $\mathbf{O}_i$ and $L_i$ into HBM, recording them as the $i$th segment in $\mathbf{O}$ and $L$.
\ENDIF
\end{algorithmic}
\end{algorithm}

In our realization of \cref{alg:flash3_wgmma_ws_only} on Hopper, we employ \verb|setmaxnreg| to manage (de)allocation, utilize TMA for transferring $\mathbf{Q}_i$ and $\{ \mathbf{K}_j, \mathbf{V}_j \}_{0 \leq j < T_c}$, and harness WGMMA for performing the GEMMs within the consumer mainloop. The SS and RS prefixes designate whether the initial operand is retrieved from shared memory or the register file, respectively. When analyzing the operational sequence of \cref{alg:flash3_wgmma_ws_only}, it is important to recognize that TMA loads are dispatched asynchronously and therefore do not block subsequent loads waiting for completion. Additionally, throughout the first $s$ cycles in the producer mainloop, no synchronization waits occur because the buffer is in the process of being populated.

\paragraph{Pingpong scheduling} The asynchronous execution of WGMMA and TMA, combined with warp-specialization, creates a possibility for overlapping the softmax computation in one warpgroup with GEMM processing in another warpgroup. This strategy is appealing because non-matmul tasks, such as special function computations, are significantly slower than matmul operations on contemporary hardware accelerators. For instance, on the H100 SXM5 GPU, the FP16 matmul performance reaches 989 TFLOPS, whereas special functions like exponential\footnote{According to the CUDA programming guide, each streaming multiprocessor (SM) can execute 16 special function operations per clock cycle. By multiplying 16 by the total number of SMs (132) and the clock speed (1830 MHz), we arrive at a throughput of 3.9 TFLOPS for special functions.} (required for softmax) achieve only 3.9 TFLOPS. In FP16 attention forward passes with a head size of 128, there are 512 times more matmul floating-point operations than exponential ones, but the hardware delivers 256 times less throughput for exponentials, resulting in the exponential step consuming about half as much time as matmul. This imbalance becomes even more pronounced with FP8, where matmul throughput doubles while the speed for exponentials remains unchanged.

Given that the exponential operation is executed by a distinct hardware component (the multi-function unit), optimal scheduling would place the exponential computation concurrent with the Tensor Cores handling the matmul tasks. To coordinate this, we introduce synchronization barriers (\texttt{bar.sync} instructions), which ensure that the GEMMs (specifically, GEMM1—$\mathbf{P} \mathbf{V}$ from a given iteration, and GEMM0—$\mathbf{Q} \mathbf{K}^\top$ from the subsequent iteration) in warpgroup 1 are executed prior to those in warpgroup 2. Consequently, while warpgroup 2 is engaged with its GEMMs, warpgroup 1 performs the softmax operation. Subsequently, the execution alternates: warpgroup 2 carries out the softmax as warpgroup 1 returns to GEMMs, hence this method is termed ``pingpong’’ scheduling. This process is visualized in~\cref{fig:pingpong_scheduling}. 

Although, in real implementations, pingpong scheduling may not precisely match the schematic illustration, empirical results indicate a notable performance gain (for example, increasing throughput from 570 TFLOPS to 620–640 TFLOPS during FP16 forward passes with a head dimension of 128 and a sequence length of 8192).

\paragraph{Attention variants} In the cases of multi-query attention~\citep{shazeer2019fast} and grouped query attention~\citep{ainslie2023gqa}, our implementation aligns with the method used in \textsc{FlashAttention-2}, modifying tensor indices so that $\mathbf{K}$ and $\mathbf{V}$ are not redundantly stored in HBM.

\subsection{Intra-warpgroup overlapping GEMMs and softmax}

Within a single warpgroup, it is possible to concurrently execute certain instructions from the softmax operation alongside selected instructions from the GEMMs. Here, we outline a method that facilitates such overlap.

In the attention algorithm, operations within the inner loop (main loop) have sequential dependencies that impede parallelization within a single iteration. For example, (local) softmax (lines \ref{code-ws:softmax_start} to \ref{code-ws:softmax_end}) relies on the output $\mathbf{S}_i^{(j)}$ of the first GEMM, while the second GEMM takes its result $\widetilde{\mathbf{P}}_i^{(j)}$ as an operand. Indeed, the wait statements in lines \ref{alg:ws_only_gemm1} and \ref{alg:ws_only_gemm2} of \cref{alg:flash3_wgmma_ws_only} serialize the execution of softmax and GEMMs. However, we can break these dependencies by pipelining across iterations through additional buffers in registers. Pursuing this idea, we propose the following two-stage\footnote{Note that the number of stages of the overlapping scheme is bounded by, but need not equal, the number $s$ of stages in the circular SMEM buffer.} GEMM-softmax pipelining algorithm:

\begin{algorithm}[H]
    \caption{\small\label{alg:flash3_wgmma}\textsc{FlashAttention-3} consumer warpgroup forward pass}
    \begin{algorithmic}[1]
      \REQUIRE  Matrices $\mathbf{Q}_i \in \mathbb{R}^{B_r \times d}$ and $\mathbf{K}, \mathbf{V} \in \mathbb{R}^{N \times d}$ in HBM, key block size $B_c$ with $T_c = \lceil \frac{N}{B_c} \rceil$.
      \STATE Dynamically adjust allocation of registers according to the number of consumer warps.
      \STATE Set up $\mathbf{O}_i = (0) \in \mathbb{R}^{B_r \times d}$ and initialize $\ell_i, m_i = (0), (-\infty) \in \mathbb{R}^{B_r}$ on-chip.
      \STATE Wait until $\mathbf{Q}_i$ and $\mathbf{K}_0$ have been loaded into shared memory.
      \STATE Execute matrix multiplication $\mathbf{S}_{\mathrm{cur}} = \mathbf{Q}_i \mathbf{K}_0^T$ via WGMMA; commit operation and synchronize.
      \STATE Release the buffer's $0$th stage allocated for $\mathbf{K}$.
      \STATE Update $m_i$, compute $\tilde{\mathbf{P}}_{\mathrm{cur}}$ and $\ell_{i}$ based on $\mathbf{S}_{\mathrm{cur}}$, and adjust $\mathbf{O}_i$ by rescaling.
      \FOR{$1 \le j < T_c - 1$}
        \STATE Await loading of $\mathbf{K}_j$ into shared memory.
        \label{code:mainloop_start}
        \STATE Calculate $\mathbf{S}_{\mathrm{next}} = \mathbf{Q}_i \mathbf{K}_j^T$ utilizing WGMMA; commit but do not synchronize. \label{code:first_wgmma}
        \STATE Await loading of $\mathbf{V}_{j-1}$ into shared memory.
        \STATE Perform the update $\mathbf{O}_i = \mathbf{O}_i + \tilde{\mathbf{P}}_{\mathrm{cur}} \mathbf{V}_{j-1}$ through WGMMA; commit without waiting. \label{code:second_wgmma}
        \STATE Wait until the WGMMA operation $\mathbf{Q}_i \mathbf{K}_j^T$ completes.
        \STATE Derive $m_{i}$, $\tilde{\mathbf{P}}_{\mathrm{next}}$, and $\ell_{i}$ utilizing $\mathbf{S}_{\mathrm{next}}$. \label{code:softmax}
        \STATE Once WGMMA operation for $\tilde{\mathbf{P}}_{\mathrm{cur}} \mathbf{V}_{j-1}$ has finished, rescale $\mathbf{O}_i$
        \STATE Free $(j\,\%\,s)$th, and $(j-1\,\%\,s)$th stage of the buffer for $\mathbf{K}$ and $\mathbf{V}$, respectively.
        \STATE Transfer $\mathbf{S}_{\mathrm{next}}$ into $\mathbf{S}_{\mathrm{cur}}$.
        \label{code:mainloop_end}
      \ENDFOR \STATE Wait for $\mathbf{V}_{T_c - 1}$ to be available in shared memory.
      \STATE Compute
      $\mathbf{O}_{i} = \mathbf{O}_{i} + \tilde{\mathbf{P}}_{\mathrm{last}} \mathbf{V}_{T_c - 1}$ employing WGMMA. Commit and synchronize.

\STATE Finalization: Adjust $\mathbf{O}_{i}$ according to the value of $m_{i}$. Determine $L_{i}$ utilizing both $m_{i}$ and $\ell_{i}$. Store $\mathbf{O}_{i}$ and $L_{i}$ in HBM, assigning them to the $i$-th segments of $\mathbf{O}$ and $L$ respectively.

\cref{alg:flash3_wgmma} serves as an alternative to the consumer portion found in \cref{alg:flash3_wgmma_ws_only}, enabling the full implementation of the \textsc{FlashAttention-3} algorithm at FP16 precision. Conceptually, WGMMA is employed here as a shorthand for asynchronous GEMM. Throughout the main processing loop (spanning lines \ref{code:mainloop_start} through \ref{code:mainloop_end}), the second invocation of the WGMMA routine in the $j$-th iteration (line \ref{code:second_wgmma}) is executed concurrently with the softmax computations associated with iteration $j+1$ (line \ref{code:softmax}).

Although the pipelined architecture discussed previously holds promise for enhanced performance, a number of pragmatic factors must be addressed:

\paragraph{Compiler reordering} The pseudocode demonstrates an ideal flow of instructions, yet in reality, the NVCC compiler frequently rearranges code to maximize optimization. This restructuring can interfere with the intended alternation between WGMMA and non-WGMMA instructions, possibly resulting in performance degradation or unpredictable outcomes. Upon examining the generated SASS code, we observe that instruction overlap occurs as anticipated (see Section~\ref{sec:2-stage-sass}).

\paragraph{Register pressure} Achieving peak performance hinges on reducing register spilling. The 2-stage pipeline model, however, necessitates supplementary registers for managing intermediate outputs and maintaining state across pipeline steps. Notably, holding an extra $\mathbf{S}_{\mathrm{next}}$ in registers increases demand by $B_r \times B_c \times \text{sizeof}(\text{float})$ for each threadblock. This surge in register utilization may constrain the possibility of larger block sizes, which themselves often require substantial register resources. Consequently, decisions on optimal configuration should be informed by empirical profiling.

\paragraph{3-stage pipelining} Building upon the prior 2-stage method, a 3-stage pipeline variant is introduced, wherein the second WGMMA operation is overlapped with the softmax computation. This approach promises further improvement in Tensor Core efficiency but, due to the additional pipeline step, increases register requirements. This further complicates the balance between increasing tile size and pipeline depth. The specifics of the 3-stage algorithm and its performance assessment are provided in ~\cref{sec:3-stage}.

\subsection{Low-precision with FP8}
\label{sec:algofp8}

\textbf{Efficiency: layout transformations.} Executing the forward pass of \textsc{FlashAttention-3} using FP8 precision introduces unique challenges related to layout compatibility that do not arise when using FP16.

To begin, it is important to recognize that the input tensors $\mathbf{Q}$, $\mathbf{K}$, and $\mathbf{V}$ are usually stored contiguously along the head dimension. However, FP8 WGMMA requires the second GEMM's $\mathbf{V}$ tiles loaded into SMEM to be contiguous with respect to the sequence length dimension in order to meet the k-major constraint. Because the TMA load operation cannot alter which dimension is contiguous, the approach must involve either (1) transposing $\mathbf{V}$ in GMEM beforehand, or (2) performing a tile-wise transpose within the kernel after transferring $\mathbf{V}$ tiles into SMEM. For option (1), one may (1a) fuse the transpose with the epilogue of an earlier operation—such as the rotary embedding—or (1b) invoke an independent preprocessing transpose kernel\footnote{A highly efficient transpose kernel can attain performance near the device’s bandwidth~\citep{colfax_cutlass_transpose_2024}.}, which swaps the strides between sequence length and head dimensions. Yet, method (1a) is challenging to implement into generic libraries, while method (1b) results in unnecessary memory overhead during inference, where memory bandwidth is a critical resource.

For FP8 \textsc{FlashAttention-3}, we select approach (2). Within the in-kernel transpose phase, we exploit the LDSM (\verb|ldmatrix|) and STSM (\verb|stmatrix|) GPU instructions, which enable a warp of threads to collaboratively transfer data between SMEM and RMEM at 128-byte segments. \footnote{The PTX specification refers to LDSM/STSM as operations that copy $8 \times 8$ blocks of 16-bit elements \cite[\S 9.7.13.4.15-16]{ptx}, though for FP8 we pack pairs of 8-bit values so these instructions are still applicable. However, the transpose versions of LDSM/STSM are unable to separate packed 8-bit values, thereby making it necessary to incorporate additional register manipulations between LDSM and STSM to successfully achieve a tile-wise transpose; specific implementation details are omitted here.} Both LDSM and STSM exhibit high register efficiency, allowing them to be deployed directly in the producer warpgroup, while supporting the rearrangement of memory layouts in tandem with copying. Additionally, starting from the second iteration, we schedule the transposition of the subsequent $\mathbf{V}$ tile so that it occurs concurrently with the execution of two WGMMAs which utilize the previous $\mathbf{V}$ tile and the active $\mathbf{K}$ tile.

Second, we note that the FP32 accumulator used in FP8 WGMMA possesses a distinct memory organization compared to the layout assumed for operand A when stored in registers, which is not the case for FP16. Illustrations in \cref{fig:rmem_accum} and \cref{fig:rmem_operand} highlight fragments of these layouts, showing the order in which values are retained in thread-specific registers. Utilizing byte permute instructions allows us to convert the format of the accumulator from the first WGMMA into one that the second WGMMA can process, aligning with the arrangement of the $\mathbf{V}$ tile resulting from the in-kernel transpose. In particular, as shown in \cref{fig:rmem_accum}, we reorder the bytes sequentially to
$$\{ \verb|d0 d1 d4 d5 d2 d3 d6 d7| \},$$
applying this pattern repeatedly for every 8-byte segment. On a logical level, this transformation permutes the columns of the $\mathbf{P}$ tile (for instance, columns $0189$ are promoted to the first four columns). To ensure WGMMA generates the correct output tile, the in-kernel transpose can be implemented to produce a corresponding permutation in the rows of the $\mathbf{V}$ tile.\footnote{The flexibility to reorder in the in-kernel transpose removes the necessity for shuffle instructions to exchange register ownership across threads, as was previously discussed in~\citep{colfax_fp8_flashattention_2024}.}

\textbf{Accuracy: block quantization and incoherent processing.} The FP8 (e4m3) format allocates 3 bits for the mantissa and 4 bits for the exponent, offering less numerical precision than either FP16 or BF16. Large-scale models frequently contain outlier entries~\citep{dettmers2208llm, sun2024massive} whose magnitudes greatly exceed the majority of the other values, complicating the quantization process. A common strategy is to apply per-tensor scaling~\citep{micikevicius2022fp8}, where each tensor receives its own scalar scaling factor (e.g., one each for $\mathbf{Q}$, $\mathbf{K}$, and $\mathbf{V}$). To enhance the accuracy of attention mechanisms in the FP8 regime, we incorporate two main strategies:

\begin{enumerate}

\item \textbf{Block quantization}: We retain one scalar for each block, partitioning each tensor $\mathbf{Q}$, $\mathbf{K}$, and $\mathbf{V}$ into blocks with dimensions $B_r \times d$ or $B_c \times d$, and quantize each block independently. This quantization step can be efficiently integrated with preceding operations—such as rotary embeddings—without causing additional performance degradation, given that these operations are limited by memory bandwidth. Because the \textsc{FlashAttention-3} algorithm is inherently block-oriented, we are able to rescale each block of $\mathbf{S}$ to reflect the block quantization, without incurring any computational overhead.

\item \textbf{Incoherent processing}: To mitigate the impact of outliers, we pre-process $\mathbf{Q}$ and $\mathbf{K}$ by multiplying them with a randomly chosen orthogonal matrix $\mathbf{M}$ before FP8 quantization. Since $\mathbf{M}$ is orthogonal, it satisfies $\mathbf{M} \mathbf{M}^\top = I$, ensuring that $(\mathbf{Q} \mathbf{M})(\mathbf{K} \mathbf{M})^\top$ remains equivalent to $\mathbf{Q} \mathbf{K}^\top$; thus, applying $\mathbf{M}$ to both $\mathbf{Q}$ and $\mathbf{K}$ does not alter the resultant attention. This randomization process effectively distributes outlier values, as every component of $\mathbf{Q} \mathbf{M}$ and $\mathbf{K} \mathbf{M}$ becomes a weighted combination of the original entries, decreasing quantization error. In practice, adopting the approach of \citet{chee2024quip} and~\citet{tseng2024quip}, we select $\mathbf{M}$ as a product of random diagonal matrices (with entries $\pm 1$) and a Hadamard matrix. This allows multiplication in $O(d \log d)$ time, a significant improvement over the standard $O(d^2)$ complexity, and this multiplication can likewise be merged with the rotary embedding without additional computational expense.
\end{enumerate}

Empirical evaluation confirms that employing these two methods reduces numerical errors by as much as 2.6$\times$, as demonstrated in \cref{sec:numerical_error}.

\section{Empirical Validation}
\label{sec:exp}

To assess both the performance and correctness of \textsc{FlashAttention-3}, we leverage CUTLASS~\citep{Thakkar_CUTLASS_2023} primitives including WGMMA and TMA abstractions in its implementation.

\begin{itemize}

\item \textbf{Benchmarking attention.} We evaluate the performance of \textsc{FlashAttention-3} at various sequence lengths, comparing its runtime to the standard PyTorch implementation, \textsc{FlashAttention-2}, the Triton version of \textsc{FlashAttention-2} (which leverages H100-specific instructions), and a vendor-provided \textsc{FlashAttention-2} implementation in cuDNN, optimized for H100 GPUs. Our results show that \textsc{FlashAttention-3} achieves up to a 2.0$\times$ speedup over \textsc{FlashAttention-2}, and is up to 1.5$\times$ faster than the Triton-based variant. Additionally, \textsc{FlashAttention-3} attains performance of up to 740 TFLOPs/s, corresponding to 75\% of the H100 GPU’s theoretical TFLOPs/s peak.
\item \textbf{Ablation study.} Our ablation analysis demonstrates that the incorporation of warp-specialization techniques and the GEMM-softmax pipelining in the algorithm underlies the performance improvements observed in \textsc{FlashAttention-3}.
\item \textbf{Accuracy of FP8 attention.} We demonstrate that the integration of block quantization and incoherent processing decreases the numerical error in FP8 \textsc{FlashAttention-3} by a factor of 2.6$\times$.

\subsection{Benchmarking Attention}
\label{subsec:benchmark_attn}

We evaluate the execution time of various attention mechanisms on an H100 80GB SXM5 GPU under several configurations (with and without a causal mask, head dimensions set to 64 or 128) using FP16 precision inputs. The performance metrics are summarized in~\cref{fig:benchmark_attn_fwd} and~\cref{fig:benchmark_attn_bwd}. Our findings indicate that \textsc{FlashAttention-3} achieves a speedup of approximately 1.5--2.0$\times$ over \textsc{FlashAttention-2} during the forward pass, and is 1.5--1.75$\times$ faster during the backward pass. In relation to a baseline attention implementation, \textsc{FlashAttention-3} attains speed improvements of up to 3--16$\times$. Notably, for sequence lengths of 1k tokens or more, \textsc{FlashAttention-3} even outperforms the cuDNN vendor library (which is closed source and optimized for H100 GPUs) in terms of throughput.

\paragraph{Benchmark settings:} We vary the sequence length as 512, 1k, ..., 16k, and set batch size so that the total number of tokens is 16k. We set the hidden dimension to 2048, and head dimension to be either 64, 128, or 256 (i.e., 32 heads, 16 heads, or 8 heads). To calculate the FLOPs of the forward pass, we use:

\begin{equation*}
  4 \cdot \text{seqlen}^2 \cdot \text{head dimension} \cdot \text{number of heads}.
\end{equation*}

Utilizing causal masking entails halving this value, as roughly 50% of the computations are performed. The backward pass FLOPs can then be determined by multiplying the forward pass FLOPs by 2.5, reflecting the ratio of matrix multiplications: 2 in the forward pass versus 5 in the backward pass (owing to recomputation).

Additionally, we evaluate the runtime performance of FP8 during the forward pass using comparable configurations. The findings for headdim 256 are presented in ~\cref{fig:benchmark_attn_fp8}, while comprehensive results are detailed in \cref{sec:benchmark_attn_fp8_full}.

\subsection{Ablation Study: 2-Stage Pipelining Experiments}
\label{sec:2-stage-pipelining-experiments}

We examine the effects of removing both the 2-stage WGMMA-softmax pipeline and warp-specialization within the context of non-causal FP16 \textsc{FlashAttention-3}, using fixed configuration parameters $\{\text{batch}, \text{seqlen}, \text{nheads}, \text{hdim}\} = \{ 4, 8448, 16, 128\}$. The findings, presented in~\cref{table:ablation_pipelining}, demonstrate that our enhancements—incorporating asynchrony via warp-specialization and enabling overlap between the GEMM and softmax operations—substantially improve throughput, boosting performance from 570 to 661 TFLOPs.

\subsection{Numerical Error Validation}
\label{sec:numerical_error}

Due to growing concerns regarding the numerical accuracy of \textsc{FlashAttention}~\citep{golden2024flash}, we conduct an evaluation of \textsc{FlashAttention-2}, \textsc{FlashAttention-3}, and a conventional attention mechanism, benchmarking each against a reference attention computation performed in FP64 precision. To emulate the presence of anomalous features and activations commonly observed in LLMs~\citep{dettmers2208llm, sun2024massive}, we initialize the components of $\mathbf{Q}, \mathbf{K}, \mathbf{V}$ using the following probabilistic scheme:

\begin{equation*}
  \mathcal{N}(0, 1) + \mathcal{N}(0, 100) \cdot \mathrm{Bernoulli}(0.001).
\end{equation*}

In this setup, every element is drawn from a normal distribution with a mean of zero and standard deviation of 1. However, an independent Gaussian noise component with standard deviation 10 is added to 0.1\% of these elements. Subsequently, we compute the root mean squared error (RMSE) as shown in~\cref{table:numerical_error}. In the FP16 format, both \textsc{FlashAttention-2} and \textsc{FlashAttention-3} show a 1.7$\times$ reduction in RMSE compared to the conventional approach, due to maintaining intermediate softmax computations in FP32. For FP8, the reference attention utilizes per-tensor scaling, retaining FP32 accumulators for matmul operations and storing intermediate softmax outputs in FP16. Owing to the use of block quantization and incoherent processing strategies, \textsc{FlashAttention-3} achieves a 2.6$\times$ improvement in accuracy over this FP8 baseline.

\section{Dicussion, Limitations, Conclusion}
\label{sec:discussion}

Utilizing \textsc{FlashAttention-3}, we have shown that advancements in programming methods and hardware capabilities—specifically asynchrony and reduced numerical precision—can significantly enhance both the efficiency and fidelity of attention computations. Our approach achieves a 1.5-2.0$\times$ increase in attention speed relative to \textsc{FlashAttention-2}, and decreases FP8 numerical error by a factor of 2.6 compared to conventional per-tensor quantization. Several aspects remain for future improvement: fine-tuning for LLM inference scenarios, adopting a persistent kernel design for the FP8 kernel,\footnote{In our performance evaluation, the FP16 version of \textsc{FlashAttention-3} benefits from both a persistent kernel and effective load balancing, whereas these features are absent from the FP8 implementation. This contributes to the relatively poorer performance of FP8 \textsc{FlashAttention-3} when dealing with short sequence lengths and causal masking, in comparison to FP8 cuDNN kernels.} and investigating the implications of low-precision attention mechanisms in large-scale model training. Although this study concentrates on Hopper GPUs, we anticipate that the underlying methods are transferable to other accelerators. Ultimately, we aim for improved, more precise attention primitives to facilitate emerging applications in long-context scenarios.

\end{document}